\documentclass[11pt]{article}
\usepackage[margin=1in]{geometry}
\usepackage{hyperref}
\hypersetup{colorlinks=true,linkcolor=black,urlcolor=black}
\pagestyle{empty}
\setlength{\parindent}{0pt}
\setlength{\parskip}{0.9em}

\begin{document}

Dear Editor,

We are pleased to submit our manuscript entitled ``Competing feedback channels drive phase transitions in collective emotion'' for consideration as a Research Article in \textit{PNAS Nexus}.

Collective emotions can remain stable for long periods and then shift abruptly on social platforms, sometimes without commensurate external shocks. Existing event-driven cascade accounts often assume stability until disturbed, leaving open the possibility that fragility accumulates endogenously. In online settings, a mixed information ecology combines stabilizing mainstream signals with amplifying user-generated feedback, yet we lack a mechanistic criterion for when amplification overwhelms stabilization and the system loses resilience.

We address this gap with a minimal mixed-feedback framework in which channel balance acts as a control parameter that sets a stability boundary. The model identifies a critical threshold beyond which the neutral state destabilizes into self-sustaining polarization, and the direction--intensity decomposition clarifies why elevated activity signals heightened fragility via depletion of the moderate buffer. Mapping the parameter landscape quantifies structural vulnerability and highlights two strongest levers---media ecology and local coupling---each shifting the stability boundary by roughly two-thirds over realistic ranges. Agent-based simulations on Erd\H{o}s--R\'enyi and Barab\'asi--Albert networks show that these signatures persist under networked interactions. In large-scale Weibo discussions during and after the COVID-19 pandemic, segments with elevated activity exhibit larger polarization jumps, and, under sufficiently dense temporal sampling, greater user-generated dominance covaries with higher polarization volatility.

We believe the manuscript aligns well with the multidisciplinary scope of \textit{PNAS Nexus}. It demonstrates how competing feedback channels give rise to phase transitions, critical slowing down, and symmetry breaking in a sociotechnical system---bridging statistical physics, computational social science, and communication studies. The unified theory-to-data evidence chain, spanning mean-field analysis, network simulations, and large-scale empirical validation, exemplifies the rigorous, cross-disciplinary approach the journal champions. Beyond theoretical interest, the framework offers actionable levers for platform governance and crisis communication.

We confirm that this manuscript has not been published elsewhere and is not under consideration by another journal. All authors have approved the submission.

Thank you for considering our submission.

Sincerely,

\textit{On behalf of all authors,}

\vspace{0.5em}

\textbf{Corresponding authors:}

Jinlin Wu\\
Urban Governance and Design Thrust\\
The Hong Kong University of Science and Technology (Guangzhou)\\
Email: jwu923@connect.hkust-gz.edu.cn

Zhihang Liu\\
Institute of Space and Earth Information Science\\
The Chinese University of Hong Kong\\
Email: zhihangliu@cuhk.edu.hk

\end{document}
